% ============================================================
%  EXAMEN TEMA 5: QUÍMICA — MATERIA, MEZCLAS Y SEPARACIONES
%  Curso Introductorio — 3er Año
%  Duración: 90 minutos
%  Temas: Método científico, laboratorio, propiedades de la
%         materia, taxonomía de la materia y separación de mezclas
%  Autor: Ing. Gabriel Astudillo
%  Nota: enunciado limpio (sin pistas ni desarrollo). Las
%  soluciones están en la "Clave de Respuestas" al final.
% ============================================================

\documentclass[12pt, letterpaper]{article}

% ── Paquetes ─────────────────────────────────────────────────

\usepackage{mathwizards-guia}
\mwguia{Examen — Tema 5: Química: Materia, Mezclas y Separaciones}{Ing. Gabriel Astudillo $\cdot$ Curso 3er Año}

\begin{document}
% ============================================================

% ── Portada / Título ─────────────────────────────────────────
\mwportada{Examen del Tema 5}{Química: Materia, Mezclas y Separaciones}{Método Científico, Propiedades de la Materia y Técnicas de Separación}

\vspace{4pt}
\begin{cuadroinfo}
  \small
  \textbf{Nombre:} \rule{0.55\linewidth}{0.4pt} \quad \textbf{Fecha:} \rule{0.15\linewidth}{0.4pt} \\[6pt]
  \textbf{Tiempo disponible:} 90 minutos \quad$\bullet$\quad \textbf{Puntaje total:} 100 puntos \\[4pt]
  \textbf{Instrucciones:}
  \begin{enumerate}[leftmargin=*]
    \item Lee cada ejercicio con calma antes de resolverlo.
    \item Justifica siempre tus respuestas; no basta con nombrar el concepto.
    \item En los cálculos, muestra el procedimiento y no olvides las unidades.
    \item Cuando diseñes procedimientos, explica qué propiedad se aprovecha en cada paso.
    \item No se permite el uso de calculadora.
    \item Escribe con letra clara y revisa tu examen antes de entregarlo.
  \end{enumerate}
\end{cuadroinfo}

\vspace{8pt}

% ============================================================
\section{El Método Científico y el Laboratorio (20 puntos)}
% ============================================================

\begin{enumerate}[leftmargin=*]
  \item \textbf{(6 pts)} Ordena correctamente las siguientes etapas del método científico:
        \[
          \text{Conclusión}, \quad \text{Observación}, \quad \text{Hipótesis}, \quad \text{Experimentación}, \quad \text{Planteamiento del problema}, \quad \text{Análisis de resultados}
        \]

  \item \textbf{(6 pts)} Clasifica cada instrumento según su categoría (sostén, uso específico, volumétrico o recipiente):
        \[
          a)\ \text{probeta graduada} \quad b)\ \text{soporte universal} \quad c)\ \text{mechero de Bunsen} \quad d)\ \text{vaso de precipitados} \quad e)\ \text{pipeta} \quad f)\ \text{gradilla}
        \]

  \item \textbf{(8 pts)} Un grupo de estudiantes quiere determinar si la temperatura del agua afecta la velocidad con que se disuelve el azúcar.
        \begin{enumerate}[label=\alph*)]
          \item Formula una hipótesis.
          \item Identifica la variable independiente, la variable dependiente y al menos dos variables controladas.
        \end{enumerate}
\end{enumerate}

\newpage

% ============================================================
\section{La Materia y sus Propiedades (25 puntos)}
% ============================================================

\begin{enumerate}[leftmargin=*]
  \item \textbf{(6 pts)} Clasifica cada propiedad como extensiva (E) o intensiva (I):
        \[
          a)\ \text{masa} \quad b)\ \text{densidad} \quad c)\ \text{volumen} \quad d)\ \text{punto de fusión} \quad e)\ \text{peso} \quad f)\ \text{color}
        \]

  \item \textbf{(6 pts)} Un objeto tiene una masa de $8\;\text{kg}$. Calcula su peso en la Tierra ($g = 9{,}8\;\text{m/s}^2$) y en la Luna ($g = 1{,}62\;\text{m/s}^2$).

  \item \textbf{(6 pts)} Un estudiante tiene dos trozos de un metal desconocido. El primero tiene masa $150\;\text{g}$ y volumen $55{,}6\;\text{cm}^3$; el segundo, masa $300\;\text{g}$ y volumen $111{,}1\;\text{cm}^3$. Calcula la densidad de cada uno y determina si son del mismo material.

  \item \textbf{(7 pts)} Un cuerpo pesa $245\;\text{N}$ en la Tierra ($g = 9{,}8\;\text{m/s}^2$).
        \begin{enumerate}[label=\alph*)]
          \item Calcula su masa.
          \item \textquestiondown Cuánto pesaría en Marte ($g_{\text{Marte}} = 3{,}72\;\text{m/s}^2$)?
          \item \textquestiondown En cuál de los dos planetas pesa más? Justifica.
        \end{enumerate}
\end{enumerate}

\newpage

% ============================================================
\section{Taxonomía de la Materia: Sustancias y Mezclas (25 puntos)}
% ============================================================

\begin{enumerate}[leftmargin=*]
  \item \textbf{(8 pts)} Clasifica cada muestra como elemento, compuesto, mezcla homogénea o mezcla heterogénea:
        \[
          a)\ \text{agua destilada} \quad b)\ \text{aire} \quad c)\ \text{oro puro} \quad d)\ \text{ensalada}
        \]
        \[
          e)\ \text{bronce} \quad f)\ \text{dióxido de carbono} \quad g)\ \text{leche} \quad h)\ \text{agua con aceite}
        \]

  \item \textbf{(8 pts)} Explica con tus palabras y con un ejemplo la diferencia entre:
        \begin{enumerate}[label=\alph*)]
          \item un \textbf{elemento} y un \textbf{compuesto};
          \item una \textbf{mezcla homogénea} y una \textbf{mezcla heterogénea}.
        \end{enumerate}

  \item \textbf{(9 pts)} Clasifica según la clasificación granulométrica (dispersión grosera, coloidal o solución):
        \[
          a)\ \text{agua de mar} \quad b)\ \text{jugo de naranja con pulpa} \quad c)\ \text{niebla} \quad d)\ \text{alcohol disuelto en agua}
        \]
        Luego explica qué es el \textbf{efecto Tyndall} y cómo se usaría para distinguir una solución de una dispersión coloidal.
\end{enumerate}

\newpage

% ============================================================
\section{Métodos de Separación de Mezclas (30 puntos)}
% ============================================================

\begin{enumerate}[leftmargin=*]
  \item \textbf{(8 pts)} Indica el método de separación más adecuado para cada mezcla y la propiedad física que se aprovecha:
        \[
          a)\ \text{arena y agua} \quad b)\ \text{agua y aceite} \quad c)\ \text{sal disuelta en agua} \quad d)\ \text{alcohol y agua}
        \]

  \item \textbf{(10 pts)} Tienes una mezcla de \textbf{limaduras de hierro, arena y sal}. Diseña un procedimiento paso a paso para separar los tres componentes, indicando en cada paso el método usado y la propiedad aprovechada.

  \item \textbf{(12 pts)} Tienes una mezcla de \textbf{arena, sal y agua}.
        \begin{enumerate}[label=\alph*)]
          \item Explica paso a paso cómo separarías los tres componentes.
          \item Si además quisieras \textbf{recuperar el agua}, \textquestiondown qué método usarías en el último paso en lugar de la evaporación? Justifica.
        \end{enumerate}
\end{enumerate}

\newpage

% ============================================================
\ifmwclaves
  \section{Clave de Respuestas}
  % ============================================================

  \subsection*{Sección 1: El Método Científico y el Laboratorio}

  \begin{enumerate}[leftmargin=*, label=\textbf{\arabic*.}]
    \item Orden correcto: Observación $\rightarrow$ Planteamiento del problema $\rightarrow$ Hipótesis $\rightarrow$ Experimentación $\rightarrow$ Análisis de resultados $\rightarrow$ Conclusión.

    \item a) volumétrico \quad b) sostén \quad c) uso específico \quad d) recipiente \quad e) volumétrico \quad f) sostén

    \item a) Hipótesis: «El azúcar se disuelve más rápido en agua caliente que en agua fría.»
          \quad b) Variable independiente: temperatura del agua. Variable dependiente: tiempo que tarda el azúcar en disolverse. Variables controladas (ejemplos): cantidad de azúcar, cantidad de agua y tipo de agitación.
  \end{enumerate}

  \newpage

  \subsection*{Sección 2: La Materia y sus Propiedades}

  \begin{enumerate}[leftmargin=*, label=\textbf{\arabic*.}]
    \item a) E \quad b) I \quad c) E \quad d) I \quad e) E \quad f) I

    \item Tierra: $W = m\,g = 8 \times 9{,}8 = 78{,}4\;\text{N}$. \quad Luna: $W = 8 \times 1{,}62 = 12{,}96\;\text{N}$.
          (La masa sigue siendo $8\;\text{kg}$ en ambos lugares.)

    \item $d_1 = \dfrac{150}{55{,}6} \approx 2{,}70\;\text{g/cm}^3$; \quad $d_2 = \dfrac{300}{111{,}1} \approx 2{,}70\;\text{g/cm}^3$.
          Las densidades coinciden, por lo que \textbf{sí son del mismo material} (probablemente aluminio).

    \item a) $m = \dfrac{W}{g} = \dfrac{245}{9{,}8} = 25\;\text{kg}$.
          \quad b) $W_{\text{Marte}} = 25 \times 3{,}72 = 93\;\text{N}$.
          \quad c) Pesa más en la \textbf{Tierra} ($245\;\text{N} > 93\;\text{N}$), porque allí la gravedad es mayor.
  \end{enumerate}

  \newpage

  \subsection*{Sección 3: Taxonomía de la Materia: Sustancias y Mezclas}

  \begin{enumerate}[leftmargin=*, label=\textbf{\arabic*.}]
    \item a) compuesto \quad b) mezcla homogénea \quad c) elemento \quad d) mezcla heterogénea \quad e) mezcla homogénea \quad f) compuesto \quad g) coloide (mezcla heterogénea a nivel microscópico) \quad h) mezcla heterogénea

    \item a) Un \textbf{elemento} no se puede descomponer en sustancias más simples por métodos químicos (ej.: hierro, Fe). Un \textbf{compuesto} está formado por dos o más elementos combinados en proporción fija y sí se puede descomponer químicamente (ej.: agua, H$_2$O).
          \\[2pt]
          b) En una \textbf{mezcla homogénea} no se distinguen sus componentes a simple vista (una sola fase; ej.: agua salada). En una \textbf{mezcla heterogénea} sí se distinguen (dos o más fases; ej.: agua con aceite).

    \item a) solución \quad b) dispersión grosera \quad c) coloidal \quad d) solución.
          El \textbf{efecto Tyndall} es la dispersión de un haz de luz al atravesar un coloide. Si se hace pasar un rayo de luz (por ejemplo, de una linterna) a través de la muestra: si el haz se ve dispersado, es un coloide; si la luz pasa sin dispersarse, es una solución verdadera.
  \end{enumerate}

  \newpage

  \subsection*{Sección 4: Métodos de Separación de Mezclas}

  \begin{enumerate}[leftmargin=*, label=\textbf{\arabic*.}]
    \item a) \textbf{Filtración}: diferencia de tamaño de partícula (el sólido queda retenido).
          \quad b) \textbf{Decantación}: diferencia de densidad (son inmiscibles).
          \quad c) \textbf{Evaporación}: diferencia de puntos de ebullición (el agua se evapora y queda la sal).
          \quad d) \textbf{Destilación}: diferencia de puntos de ebullición (se recuperan ambos líquidos).

    \item Procedimiento:
          \begin{enumerate}[label=\arabic*.]
            \item \textbf{Imantación}: se acerca un imán para extraer las \textbf{limaduras de hierro}. Propiedad: magnetismo.
            \item \textbf{Disolución + filtración}: se añade agua a la mezcla restante; la sal se disuelve y la arena no. Se filtra para separar la \textbf{arena} (queda en el filtro). Propiedad: solubilidad y tamaño de partícula.
            \item \textbf{Evaporación}: se calienta el agua salada; el agua se evapora y queda la \textbf{sal} en el recipiente. Propiedad: diferencia de puntos de ebullición.
          \end{enumerate}

    \item a) Separación de arena, sal y agua:
          \begin{enumerate}[label=\arabic*.]
            \item \textbf{Filtración}: se separa la \textbf{arena} (sólido insoluble) del agua salada. Propiedad: tamaño de partícula.
            \item \textbf{Evaporación}: se evapora el agua y queda la \textbf{sal}. Propiedad: diferencia de puntos de ebullición.
          \end{enumerate}
          \textbf{b)} Para \textbf{recuperar el agua} se usa la \textbf{destilación} en lugar de la evaporación: en la destilación el vapor de agua se conduce a un refrigerante, donde se condensa y se recoge como agua líquida, separándola también de la sal.
  \end{enumerate}

\fi
\end{document}
